\documentclass[11pt,aspectratio=169]{beamer}

% ── Кодировка и языки ──────────────────────────────────────────────
\usepackage[T2A]{fontenc}
\usepackage[utf8]{inputenc}
\usepackage[english,russian]{babel}

% ── Тема оформления ─────────────────────────────────────────────────
\usetheme{metropolis}
\usefonttheme{professionalfonts}
\setbeamerfont{frametitle}{size=\large,series=\bfseries}
\setbeamerfont{block title}{size=\small,series=\bfseries}

% ── Пакеты ──────────────────────────────────────────────────────────
\usepackage{booktabs}
\usepackage{graphicx}
\usepackage{tikz}
\usepackage{pgfplots}
\pgfplotsset{compat=1.18}
\usepackage{xcolor}
\usepackage{colortbl}

% ── Цвета ───────────────────────────────────────────────────────────
\definecolor{myblue}{RGB}{0,102,204}
\definecolor{mygreen}{RGB}{0,153,76}
\definecolor{myred}{RGB}{204,0,51}
\definecolor{myorange}{RGB}{255,128,0}
\definecolor{lightgray}{RGB}{240,240,240}

\setbeamercolor{alerted text}{fg=myred}
\setbeamercolor{example text}{fg=mygreen}

% ── Титульная информация ────────────────────────────────────────────
\title{Применение depthwise separable convolutions для ускорения инференса CNN на CPU}
\subtitle{Курсовая работа}
\author{Станиславчук Сергей}
\date{3 июня 2026}
\institute{}

\begin{document}

% ══════════════════════════════════════════════════════════════════════
\maketitle
% ══════════════════════════════════════════════════════════════════════

% ── Слайд 2: Введение ───────────────────────────────────────────────
\begin{frame}{Введение}
  \begin{block}{Проблема}
    Глубокие CNN (ResNet-18, ResNet-50, \dots) достигают высокой точности,
    но содержат \textbf{десятки миллионов параметров} и требуют больших
    вычислительных затрат. Это ограничивает их применение на мобильных
    и~встраиваемых устройствах.
  \end{block}

  \pause

  \begin{block}{Решение}
    Замена \textit{стандартных} свёрток на \textit{depthwise-разделимые}
    (в стиле MobileNet) позволяет резко сократить число параметров
    и~ускорить инференс при минимальной потере точности.
  \end{block}

  \pause

  \begin{block}{Цель работы}
    Реализовать и сравнить ResNet-18 со стандартными (1) и~depthwise-разделимыми (2)
    свёртками на датасете CIFAR-10 по трём метрикам:
    \begin{itemize}
      \item число параметров и~размер модели
      \item задержка инференса (latency) на CPU
      \item точность классификации (test accuracy)
    \end{itemize}
  \end{block}
\end{frame}

% ── Слайд 3: Стандартная свёртка ────────────────────────────────────
\begin{frame}{Стандартная свёртка}
  \centering
  \begin{tikzpicture}[scale=0.85, every node/.style={transform shape}]
    \def\s{1.0}
    \def\ks{0.9}

    % Input feature map
    \fill[myblue!20] (0,0) rectangle (3.5*\ks,3.5*\ks);
    \draw[myblue, very thick] (0,0) rectangle (3.5*\ks,3.5*\ks);
    \node[myblue] at (1.75*\ks,3.8*\ks) {\(C_{\text{in}}\) каналов};

    % Arrow with kernel
    \draw[->, very thick] (4.0*\ks, 1.75*\ks) -- (5.5*\ks, 1.75*\ks);
    \node at (4.75*\ks, 2.5*\ks) {\small \(k \times k\)};
    \node at (4.75*\ks, 1.0*\ks) {\footnotesize \(C_{\text{in}} \cdot C_{\text{out}}\)};

    % Output
    \fill[mygreen!20] (6.0*\ks,0.5*\ks) rectangle (9.5*\ks,3.0*\ks);
    \draw[mygreen, very thick] (6.0*\ks,0.5*\ks) rectangle (9.5*\ks,3.0*\ks);
    \node[mygreen] at (7.75*\ks,3.3*\ks) {\(C_{\text{out}}\) каналов};

    % Formula
    \node[anchor=west, align=left] at (0, -1.5) {
      \textbf{Параметры:} \(k^2 \cdot C_{\text{in}} \cdot C_{\text{out}}\)\\
      \textbf{Вычисления:} \(k^2 \cdot C_{\text{in}} \cdot C_{\text{out}} \cdot H \cdot W\)
    };
  \end{tikzpicture}
\end{frame}

% ── Слайд 4: Depthwise-разделимая свёртка ───────────────────────────
\begin{frame}{Depthwise-разделимая свёртка}
  \centering
  \begin{tikzpicture}[scale=0.75, every node/.style={transform shape}]
    \def\ks{0.8}

    % Stage 1: Depthwise
    \fill[myblue!20] (0,0) rectangle (3*\ks,3*\ks);
    \draw[myblue, very thick] (0,0) rectangle (3*\ks,3*\ks);
    \node[myblue] at (1.5*\ks,3.3*\ks) {\footnotesize \(C_{\text{in}}\)};

    \draw[->, very thick] (3.5*\ks, 1.5*\ks) -- (5.0*\ks, 1.5*\ks);
    \node at (4.25*\ks, 2.5*\ks) {\footnotesize depthwise};
    \node at (4.25*\ks, 0.5*\ks) {\footnotesize \(k \times k\)};

    \fill[myorange!20] (5.5*\ks,0) rectangle (8.5*\ks,3*\ks);
    \draw[myorange, very thick] (5.5*\ks,0) rectangle (8.5*\ks,3*\ks);
    \node[myorange] at (7.0*\ks,3.3*\ks) {\footnotesize \(C_{\text{in}}\)};

    % Stage 2: Pointwise
    \draw[->, very thick] (9.0*\ks, 1.5*\ks) -- (10.5*\ks, 1.5*\ks);
    \node at (9.75*\ks, 2.5*\ks) {\footnotesize pointwise};
    \node at (9.75*\ks, 0.5*\ks) {\footnotesize \(1 \times 1\)};

    \fill[mygreen!20] (11.0*\ks,0.3*\ks) rectangle (14.0*\ks,2.7*\ks);
    \draw[mygreen, very thick] (11.0*\ks,0.3*\ks) rectangle (14.0*\ks,2.7*\ks);
    \node[mygreen] at (12.5*\ks,3.0*\ks) {\footnotesize \(C_{\text{out}}\)};

    % Formulas
    \node[anchor=west, align=left] at (0, -1.8) {
      \textbf{Depthwise:} \;\; \(k^2 \cdot C_{\text{in}}\) \hspace{1cm} (каждый канал — отдельно)\\
      \textbf{Pointwise:} \; \(C_{\text{in}} \cdot C_{\text{out}}\) \hspace{0.8cm} (смешивание каналов)\\[4pt]
      \textbf{Всего:} \(k^2 \cdot C_{\text{in}} + C_{\text{in}} \cdot C_{\text{out}}\)
    };
  \end{tikzpicture}

  \pause
  \begin{block}{Сравнение по параметрам (\(k = 3\))}
    \[
      \frac{\text{DW-Sep}}{\text{Standard}}
      = \frac{9 C_{\text{in}} + C_{\text{in}} C_{\text{out}}}
             {9 C_{\text{in}} C_{\text{out}}}
      \approx \frac{1}{C_{\text{out}}} + \frac{1}{9}
    \]
    При типичном \(C_{\text{out}} \ge 64\) — сокращение в~\textbf{5--10 раз}.
  \end{block}
\end{frame}

% ── Слайд 5: Архитектура ResNet-18 для CIFAR-10 ─────────────────────
\begin{frame}{Архитектура ResNet-18 (CIFAR-10)}
  \begin{center}
    \begin{tabular}{lll}
      \toprule
      \textbf{Слой}          & \textbf{Выход}       & \textbf{Блоки}            \\
      \midrule
      conv1                  & \(32 \times 32\)    & \(3 \times 3,\; 64\)      \\
      \addlinespace
      layer1                 & \(32 \times 32\)    & \(\begin{bmatrix}
                                                        3 \times 3,\; 64 \\
                                                        3 \times 3,\; 64
                                                      \end{bmatrix} \times 2\) \\
      \addlinespace
      layer2                 & \(16 \times 16\)    & \(\begin{bmatrix}
                                                        3 \times 3,\; 128 \\
                                                        3 \times 3,\; 128
                                                      \end{bmatrix} \times 2\) \\
      \addlinespace
      layer3                 & \(8 \times 8\)      & \(\begin{bmatrix}
                                                        3 \times 3,\; 256 \\
                                                        3 \times 3,\; 256
                                                      \end{bmatrix} \times 2\) \\
      \addlinespace
      layer4                 & \(4 \times 4\)      & \(\begin{bmatrix}
                                                        3 \times 3,\; 512 \\
                                                        3 \times 3,\; 512
                                                      \end{bmatrix} \times 2\) \\
      \addlinespace
      avg-pool + fc          & \(1 \times 1\)      & 10 классов                \\
      \bottomrule
    \end{tabular}
  \end{center}

  \begin{itemize}
    \item Схема для CIFAR-10: без начального пулинга (Identity вместо MaxPool).
    \item Все \(3 \times 3\) и \(1 \times 1\) свёртки заменяются на DepthwiseSeparableConv2d
          во второй модели.
  \end{itemize}
\end{frame}

% ── Слайд 7: Датасет CIFAR-10 ───────────────────────────────────────
\begin{frame}{Датасет: CIFAR-10}
  \begin{columns}[T]
    \begin{column}{0.48\textwidth}
      \begin{block}{Характеристики}
        \begin{itemize}
          \item \textbf{10 классов}: самолёт, авто, птица, кот, олень, собака,
                лягушка, лошадь, корабль, грузовик.
          \item \textbf{60 000} цветных изображений \(32 \times 32\).
          \item 50 000 — обучение, 10 000 — тест.
        \end{itemize}
      \end{block}
    \end{column}
    \begin{column}{0.48\textwidth}
      \begin{block}{Аугментации}
        \begin{itemize}
          \item RandomCrop(32, padding=4)
          \item RandomHorizontalFlip
          \item Нормализация:
                \(\mu = (0.4914, 0.4822, 0.4465)\)
                \(\sigma = (0.2470, 0.2435, 0.2616)\)
        \end{itemize}
      \end{block}
    \end{column}
  \end{columns}
\end{frame}

% ── Слайд 8: Методология обучения ───────────────────────────────────
\begin{frame}{Методология обучения}
  \begin{columns}[T]
    \begin{column}{0.45\textwidth}
      \begin{block}{Гиперпараметры}
        \begin{itemize}
          \item Оптимизатор: SGD, момент 0.9
          \item Weight decay: \(5 \cdot 10^{-4}\)
          \item Начальный lr: 0.1
          \item MultiStepLR с \(\gamma = 0.1\)
          \item Этапы: на 50\% и 75\% эпох
          \item Batch size: 128
        \end{itemize}
      \end{block}
    \end{column}
    \begin{column}{0.55\textwidth}
      \begin{block}{Эксперимент}
        \begin{itemize}
          \item Устройство: CPU (Intel Core)
          \item Эпохи: 20
          \item Потоки CPU: 8
          \item Seed: 42 (воспроизводимость)
        \end{itemize}
        \vspace{4pt}
        \alert{Обе модели обучались в идентичных условиях} — те же данные,
        тот же seed, те же гиперпараметры.
      \end{block}
    \end{column}
  \end{columns}
\end{frame}

% ── Слайд 9: График обучения ────────────────────────────────────────
\begin{frame}{Динамика обучения}
  \centering
  \begin{tikzpicture}
    \begin{axis}[
      width=0.75\textwidth,
      height=0.55\textheight,
      xlabel={Эпоха},
      ylabel={Test Accuracy (\%)},
      xmin=1, xmax=20,
      ymin=35, ymax=95,
      grid=major,
      legend pos=south east,
    ]
      \addplot[myblue, very thick, mark=*] coordinates {
        (1,40.75)(2,56.83)(3,64.61)(4,74.26)(5,75.29)
        (6,76.69)(7,77.41)(8,77.49)(9,78.19)(10,79.51)
        (11,88.86)(12,89.48)(13,89.84)(14,90.06)(15,89.83)
        (16,90.73)(17,91.01)(18,91.01)(19,91.08)(20,91.09)
      };
      \addlegendentry{ResNet-18 (standard)};

      \addplot[mygreen, very thick, mark=square*] coordinates {
        (1,42.81)(2,48.58)(3,59.80)(4,59.78)(5,57.99)
        (6,71.62)(7,71.88)(8,71.69)(9,73.24)(10,71.87)
        (11,87.49)(12,87.91)(13,88.29)(14,88.80)(15,88.12)
        (16,89.36)(17,89.62)(18,89.47)(19,89.64)(20,89.81)
      };
      \addlegendentry{ResNet-18 (DW-Sep)};
    \end{axis}
  \end{tikzpicture}

  \begin{itemize}
    \item После снижения lr на 10-й эпохе — резкий скачок точности у обеих моделей.
    \item Separable-версия обучается чуть медленнее, но выходит на~плато рядом со~standard.
  \end{itemize}
\end{frame}

% ── Слайд 10: Параметры и размер ────────────────────────────────────
\begin{frame}{Сравнение: параметры и размер модели}
  \centering
  \begin{columns}[T]
    \begin{column}{0.48\textwidth}
      \begin{tikzpicture}
        \begin{axis}[
          width=\textwidth,
          height=0.45\textheight,
          ybar,
          bar width=20pt,
          ylabel={Параметры (млн)},
          symbolic x coords={Standard, DW-Sep},
          xtick=data,
          enlarge x limits=0.3,
          nodes near coords,
          nodes near coords align={vertical},
        ]
          \addplot[fill=myblue] coordinates {(Standard,11.17) (DW-Sep,1.44)};
        \end{axis}
      \end{tikzpicture}
    \end{column}
    \begin{column}{0.48\textwidth}
      \begin{tikzpicture}
        \begin{axis}[
          width=\textwidth,
          height=0.45\textheight,
          ybar,
          bar width=20pt,
          ylabel={Размер (MB)},
          symbolic x coords={Standard, DW-Sep},
          xtick=data,
          enlarge x limits=0.3,
          nodes near coords,
          nodes near coords align={vertical},
        ]
          \addplot[fill=mygreen] coordinates {(Standard,42.70) (DW-Sep,5.57)};
        \end{axis}
      \end{tikzpicture}
    \end{column}
  \end{columns}

  \pause
  \begin{block}{Ключевой результат}
    \centering
    \large
    \textbf{Сокращение в \(\mathbf{7.67\times}\)} по размеру модели
    \quad и \quad
    \textbf{в \(\mathbf{7.76\times}\)} по числу параметров.
  \end{block}
\end{frame}

% ── Слайд 11: Задержка инференса ────────────────────────────────────
\begin{frame}{Сравнение: задержка инференса (CPU)}
  \centering
  \begin{tikzpicture}
    \begin{axis}[
      width=0.65\textwidth,
      height=0.48\textheight,
      ybar,
      bar width=25pt,
      ylabel={Latency (ms)},
      symbolic x coords={Standard, DW-Sep},
      xtick=data,
      enlarge x limits=0.3,
      nodes near coords,
      nodes near coords align={vertical},
      ymin=0,
    ]
      \addplot[fill=myorange] coordinates {(Standard,27.53) (DW-Sep,6.93)};
    \end{axis}
  \end{tikzpicture}

  \begin{block}{Детали бенчмарка}
    \begin{itemize}
      \item Устройство: CPU (8 потоков), batch = 1, размер входа \(3 \times 32 \times 32\).
      \item Прогрев: 50 итераций, замеры: 300 итераций.
      \item Стандартная модель: \(27.53 \pm 2.33\) мс.
      \item Separable-модель: \alert{\(6.93 \pm 0.16\)} мс.
    \end{itemize}
    \textbf{Ускорение в \(\mathbf{3.97\times}\)} на CPU-инференсе.
  \end{block}

\end{frame}

% ── Слайд 12: Точность ──────────────────────────────────────────────
\begin{frame}{Сравнение: точность классификации}
  \centering
  \begin{tikzpicture}
    \begin{axis}[
      width=0.58\textwidth,
      height=0.42\textheight,
      ybar,
      bar width=22pt,
      ylabel={Test Accuracy (\%)},
      symbolic x coords={Standard, DW-Sep},
      xtick=data,
      enlarge x limits=0.3,
      ymin=85, ymax=93,
      nodes near coords,
      nodes near coords align={vertical},
    ]
      \addplot[fill=myblue] coordinates {(Standard,91.09) (DW-Sep,89.81)};
    \end{axis}
  \end{tikzpicture}

  \begin{block}{Результат}
    \begin{itemize}
      \item Standard: \textbf{91.09\%}
      \item DW-Separable: \textbf{89.81\%}
      \item Потеря точности: всего \(\mathbf{1.28}\) процентных пункта.
    \end{itemize}
  \end{block}

  \pause
  \begin{alertblock}{Вывод}
    Потеря \(\approx 1.3\) п.п. точности при сокращении модели в~\textbf{7.7 раз}
    и~ускорении в~\textbf{4 раза} — \alert{отличный компромисс}.
  \end{alertblock}
\end{frame}

% ── Слайд 13: Итоговая таблица ──────────────────────────────────────
\begin{frame}{Сводная таблица результатов}
  \centering
  \begin{tabular}{lrrrr}
    \toprule
    \textbf{Модель} & \textbf{Параметры} & \textbf{Размер} & \textbf{Latency} & \textbf{Точность} \\
                     &                     & \textbf{(MB)}   & \textbf{(ms)}    & \textbf{(\%)} \\
    \midrule
    ResNet-18 (standard)
        & 11\,173\,962 & 42.70 & \(27.53 \pm 2.33\) & 91.09 \\
    \addlinespace
    ResNet-18 (DW-Sep)
        & 1\,438\,565  & 5.57  & \(6.93 \pm 0.16\)  & 89.81 \\
    \bottomrule
  \end{tabular}

  \vspace{12pt}

  \begin{block}{Производные метрики}
    \begin{description}
      \item[Latency speedup] \(\mathbf{3.97\times}\) — separable-модель быстрее
      \item[Size ratio] \(\mathbf{7.67\times}\) — separable-модель меньше
      \item[Accuracy \(\Delta\)] \(+1.28\) п.п. — standard точнее
    \end{description}
  \end{block}
\end{frame}

% ── Слайд 14: Анализ ────────────────────────────────────────────────
\begin{frame}{Анализ результатов}
  \begin{block}{Почему такое большое сжатие?}
    \begin{itemize}
      \item В глубоких слоях (\(C = 256, 512\)) стандартная свёртка
            содержит \(9 \cdot 256 \cdot 512 \approx 1.18\) млн параметров.
      \item Depthwise-разделимая: \(9 \cdot 256 + 256 \cdot 512 \approx 0.13\) млн —
            в~\textbf{9 раз} меньше.
      \item Эффект накапливается через 17 свёрточных слоёв ResNet-18.
    \end{itemize}
  \end{block}

  \pause

  \begin{block}{Почему точность почти не падает?}
    \begin{itemize}
      \item Depthwise-свёртка сохраняет пространственную структуру (те же \(k \times k\)).
      \item Pointwise \(1 \times 1\) даёт полное смешивание каналов.
      \item Skip-соединения (ResNet) компенсируют снижение ёмкости.
    \end{itemize}
  \end{block}
\end{frame}

% ── Слайд 15: Обучение: время ───────────────────────────────────────
\begin{frame}{Время обучения}
  \centering
  \begin{tikzpicture}
    \begin{axis}[
      width=0.55\textwidth,
      height=0.42\textheight,
      ybar,
      bar width=22pt,
      ylabel={Время обучения (с)},
      symbolic x coords={Standard, DW-Sep},
      xtick=data,
      enlarge x limits=0.3,
      nodes near coords,
      nodes near coords align={vertical},
      ymin=0,
    ]
      \addplot[fill=myred!60] coordinates {(Standard,1010) (DW-Sep,639)};
    \end{axis}
  \end{tikzpicture}

  \begin{itemize}
    \item Стандартная модель: \(\approx 1010\) с (20 эпох, 8 CPU-потоков).
    \item Separable-модель: \(\approx 639\) с.
    \item Обучение separable-версии на \textbf{37\% быстрее} — меньше параметров
          \(\Rightarrow\) меньше градиентов \(\Rightarrow\) меньше вычислений.
  \end{itemize}

  \pause
\end{frame}

% ── Слайд 16: Применение ────────────────────────────────────────────
\begin{frame}{Где это применимо?}
  \begin{columns}[T]
    \begin{column}{0.48\textwidth}
      \centering
      \textbf{Мобильные устройства}
      \vspace{6pt}

      \begin{tikzpicture}[scale=0.7]
        \fill[lightgray] (0,0) rectangle (3,5);
        \fill[myblue!15] (0.3,0.3) rectangle (2.7,4.7);
        \node at (1.5,2.5) {\textbf{5.6 MB}};
        \node at (1.5,1.5) {\(\ll 43\) MB};
      \end{tikzpicture}

      \begin{itemize}
        \item Классификация изображений
        \item Распознавание объектов
        \item On-device ML
      \end{itemize}
    \end{column}
    \begin{column}{0.48\textwidth}
      \centering
      \textbf{Встраиваемые системы}
      \vspace{6pt}

      \begin{tikzpicture}[scale=0.7]
        \draw[fill=lightgray] (0,0) rectangle (3,2);
        \node at (1.5,1) {Raspberry Pi};
      \end{tikzpicture}

      \begin{itemize}
        \item Ограниченная память
        \item Нет GPU
        \item Низкое энергопотребление
      \end{itemize}
    \end{column}
  \end{columns}

  \vspace{8pt}
  \begin{alertblock}{}
    \centering
    Depthwise-разделимые свёртки — основа MobileNet, EfficientNet, Xception и~других
    эффективных архитектур.
  \end{alertblock}
\end{frame}

% ── Слайд 17: Заключение ────────────────────────────────────────────
\begin{frame}{Заключение}
  \begin{enumerate}
    \item Реализованы и~обучены две версии ResNet-18 для CIFAR-10:
          \textbf{стандартная} и~\textbf{depthwise-разделимая}.
    \item Depthwise-разделимая версия:
          \begin{itemize}
            \item в \(7.7\times\) меньше параметров,
            \item в \(7.7\times\) меньше занимает на диске,
            \item в \(4.0\times\) быстрее на CPU-инференсе,
            \item теряет лишь \(1.28\) п.п. точности (89.81\% против 91.09\%).
          \end{itemize}
    \item Подход оправдан для сценариев с ограниченными ресурсами
          (мобильные и встраиваемые устройства).
  \end{enumerate}

  \vspace{10pt}
  \begin{center}
    \large\textbf{Спасибо за внимание!}
  \end{center}
\end{frame}

% ══════════════════════════════════════════════════════════════════════
\end{document}
% ══════════════════════════════════════════════════════════════════════
